\documentclass[../../main.tex]{subfiles}

\graphicspath{{../../image/}{../../figure/}}

\begin{document}

\section{伴随平衡律}\label{section:ch01-伴随平衡律}

考虑 $\mathbb{R}^{k}$ 的开子集 $\mathscr{X}$ 上由场方程\cref{align:ch01-system-divA} 与本构关系\cref{align:ch01-constitutive} 联立所得的平衡律方程组\cref{align:ch01-closed-system}。定义在 $\mathcal{O}\times\mathscr{X}$ 上、取值于 $\mathbb{M}^{1\times k}$ 的光滑函数 $Q$ 称为 $G$ 的一个\emph{伴随}（companion），若存在定义在 $\mathcal{O}\times\mathscr{X}$ 上、取值于 $\mathbb{R}^{n}$ 的光滑函数 $B$，使得对一切 $U\in\mathcal{O}$ 与 $X\in\mathscr{X}$，
\begin{align}
\mathsf{D}Q_{\alpha}(U,X)=B(U,X)^{\top}\mathsf{D}G_{\alpha}(U,X),\qquad \alpha=1,\cdots,k.\label{align:ch01-companion}
\end{align}

\cref{align:ch01-companion} 的意义在于如下观察：平衡律方程组\cref{align:ch01-closed-system} 的任一古典解 $U$ 自动也是\emph{伴随平衡律}
\begin{align}
\operatorname{div} Q(U(X),X)=h(U(X),X),\label{align:ch01-companion-law}
\end{align}
的（古典）解，其中
\begin{align}
h(U,X)=B(U,X)^{\top}\Pi(U,X)+\nabla\cdot Q(U,X)-B(U,X)^{\top}\nabla\cdot G(U,X).\label{align:ch01-companion-source}
\end{align}
在\cref{align:ch01-companion-source} 中，$\nabla\cdot$ 表示固定 $U$ 时关于 $X$ 的散度——区别于 $\operatorname{div}$，后者把 $U$ 视为 $X$ 的函数。

人们通过识别使\cref{align:ch01-companion} 右端成为某 $U$ 函数之梯度的积分因子 $B$，来确定给定平衡律方程组\cref{align:ch01-closed-system} 的伴随平衡律\cref{align:ch01-companion-law}。相关可积性条件为
\begin{align}
\mathsf{D}B(U,X)^{\top}\mathsf{D}G_{\alpha}(U,X)=\mathsf{D}G_{\alpha}(U,X)^{\top}\mathsf{D}B(U,X),\qquad \alpha=1,\cdots,k,\label{align:ch01-integrability}
\end{align}
对一切 $U\in\mathcal{O}$ 与 $X\in\mathscr{X}$ 成立。显然，取任何不依赖于 $U$ 的 $B$ 即可满足\cref{align:ch01-integrability}；但此时所得伴随平衡律\cref{align:ch01-companion-law} 只是原方程组\cref{align:ch01-closed-system} 各方程的平凡线性组合。对于随 $U$ 变化的非平凡 $B$，\cref{align:ch01-integrability} 对 $B$ 的 $n$ 个未知分量施加了 $\frac{1}{2}n(n-1)k$ 个条件。因此当 $n=1$、$k$ 任意时可取任一（标量值）函数 $B$；当 $n=k=2$ 时，\cref{align:ch01-integrability} 化为两个未知量的两个方程，由此大概可确定一族 $B$。但在所有其他情形，\cref{align:ch01-integrability} 形式上超定，一般不应期望存在非平凡伴随平衡律。尽管如此，正如我们将在%

\cref{chapter:ch03-双曲平衡律方程组} 中看到的，连续介质物理的平衡律方程组天然配有伴随平衡律。

平衡律方程组\cref{align:ch01-closed-system} 称为\emph{对称的}，若对任意 $U\in\mathcal{O}$ 与 $X\in\mathscr{X}$，$n\times n$ 矩阵 $\mathsf{D}G_{\alpha}(U,X)$（$\alpha=1,\cdots,k$）均对称；设 $\mathcal{O}$ 单连通，则
\begin{align}
G(U,X)^{\top}=\mathsf{D}\Gamma(U,X)^{\top},\label{align:ch01-symmetric}
\end{align}
其中 $\Gamma$ 是某定义在 $\mathcal{O}\times\mathscr{X}$ 上、取值于 $\mathbb{M}^{1\times k}$ 的光滑函数。此时可取 $B(U,X)\equiv U$ 满足\cref{align:ch01-integrability}，从而诱导伴随
\begin{align}
Q(U,X)=U^{\top}G(U,X)-\Gamma(U,X).\label{align:ch01-Q-legendre}
\end{align}

反之，若对某 $B$ 有\cref{align:ch01-companion} 成立，且对每个固定 $X\in\mathscr{X}$，映射 $B(\cdot,X)$ 将 $\mathcal{O}$ 微分同胚地映到 $\mathbb{R}^{n}$ 的某开子集 $\mathcal{O}^{*}$，则状态向量的变换 $U^{*}=B(U,X)$ 把\cref{align:ch01-closed-system} 化为等价的平衡律方程组
\begin{align}
\operatorname{div} G^{*}(U^{*}(X),X)=\Pi^{*}(U^{*}(X),X),\label{align:ch01-transformed}
\end{align}
其中
\begin{align}
G^{*}(U^{*},X)=G(B^{-1}(U^{*},X),X),\qquad
\Pi^{*}(U^{*},X)=\Pi(B^{-1}(U^{*},X),X),\label{align:ch01-G-star}
\end{align}
它是对称的。事实上，令
\begin{align}
Q^{*}(U^{*},X)=Q(B^{-1}(U^{*},X),X),\label{align:ch01-Q-star}
\end{align}
\begin{align}
\Gamma^{*}(U^{*},X)=U^{*\top}G^{*}(U^{*},X)-Q^{*}(U^{*},X),\label{align:ch01-Gamma-star}
\end{align}
则由\cref{align:ch01-companion} 容易得到
\begin{align}
G^{*}(U^{*},X)^{\top}=\mathsf{D}\Gamma^{*}(U^{*},X)^{\top}.\label{align:ch01-G-star-symmetric}
\end{align}

于是我们证明了：一条平衡律方程组配有非平凡伴随平衡律，当且仅当它是\emph{可对称化的}。

我们将看到，伴随平衡律的存在对物理中出现的平衡律方程组理论有重大意义。通常为简化分析，必须作出简化物理假设以截断平衡律方程组，同时按比例缩减状态向量的维数。这类截断不能任意进行而不破坏方程组的数学结构——而这种结构恰与其物理相关性相伴。对于正则截断，必须在方程组的对称形式\cref{align:ch01-transformed} 上操作（或至少按它来思考），并遵守如下规则：删去第 $i$ 条平衡律应同时“冻结”（即赋予固定值）特殊状态向量 $U^{*}$ 的第 $i$ 个分量 $U^{*i}$。这样所得截断方程组仍保持对称，并继承伴随
\begin{align}
\hat{Q}=Q^{*}-\sum U^{*i}G^{*i},\label{align:ch01-truncated-companion}
\end{align}
其中求和遍历所有已删去第 $i$ 条平衡律且 $U^{*i}$ 已被冻结的 $i$，$G^{*i}$ 表示 $G^{*}$ 的第 $i$ 行向量。

尽管有\cref{align:ch01-companion}，但与古典解的行为不同，\cref{align:ch01-closed-system} 的弱解未必满足\cref{align:ch01-companion-law}。然而平衡律方程组理论的信条之一是：\emph{容许弱解}至少应对指定的一族伴随满足不等式
\begin{align}
\operatorname{div} Q(U(X),X)\leqslant h(U(X),X),\label{align:ch01-admissibility}
\end{align}
在分布意义下。把这一公设与热力学第二定律联系起来，并研究它对弱解稳定性的意涵，是本书的主要目标之一。

注意到，在分布意义下成立的不等式\cref{align:ch01-admissibility}，总可通过从右端减去某非负测度 $\mathscr{M}$ 而化为等式：
\begin{align}
\operatorname{div} Q(U(X),X)=h(U(X),X)-\mathscr{M},\label{align:ch01-admissibility-equality}
\end{align}
从而由\cref{theorem:ch01-一般迹}，可将其实现为一条平衡律的场方程。

\end{document}
